\section{A new source invariant and two regions of unreachable code}
\label{sec:source-product-invariant}

The following result concerns the exact 6,455-leaf source program underlying
\texttt{zfc\_composed237\_old.tm}.  It is an ordinary inductive proof supported
by a finite certificate, and it does not itself reduce the binary state count.
It supplies two precise regions in which source code may be changed without
affecting the initialized proof-search computation.

The source list has SHA-256
\begin{quote}\small\ttfamily
7fd152d1f5371d2a561d4ad1870d241683\allowbreak
b010ce7b7a27399b95363a4537cfb6.
\end{quote}
Its ordered instruction indices start at zero.  Registers are logical
registers, before the inherited physical permutation.  The initial register
file is
\[
(3,0,0,0,3,0,0,0,0).
\]
The finite graph interface is the inherited ordered-leaf graph: an increment
advances the program counter, a positive conditional decrement decreases its
register and advances, and a zero conditional decrement executes a level-one
break.  The graph extraction checks the required zero-padded path convention.

\subsection{A finite abstraction with an unbounded meaning}

Use the eighteen classes
\[
\{0\},\{1\},\ldots,\{15\},\qquad
 E=\{16,18,20,\ldots\},\qquad
 O=\{17,19,21,\ldots\}.
\]
An abstract register value is any union of these classes.  An abstract register
file is the Cartesian product of nine such unions.  Different program points
have different abstract register files; an unreached point is assigned the
empty set.  Class unions are encoded as eighteen-bit nonnegative integers.

Increment is exact on the singleton classes below 15, maps 15 into $E$, maps
$E$ into $O$, and maps $O$ into $E$.  Positive decrement is exact on positive
singletons, and its two large-class images are
\[
\operatorname{pred}(E)=\{15\}\cup O,
\qquad
\operatorname{pred}(O)=E.
\]
Transfer across a zero decrement branch restricts its tested register to
$\{0\}$; transfer across a positive branch first excludes zero and then applies
positive decrement.  Jumps preserve all register values.  Every abstract image described here contains the concrete image.  The
displayed predecessor images of $E$ and $O$ are exact; mapping 15 into $E$
and incrementing $O$ can overapproximate the concrete image.  Taking Cartesian
products can additionally discard relations between registers, but it only
enlarges the possible set of files.

Starting with the exact initialized seed at the first decoded instruction,
the replay iterates these transfers and componentwise union until a fixed
point is reached.  The lattice is finite, so this process terminates without
a bound on the source computation, register values, or the number of proof
search iterations.  For the threshold-16 run, 273,541 strict point updates
produce the retained certificate.  A separate pass checks every applicable
edge against the final annotation.

\paragraph{Inductive invariant theorem.}
Every initialized source configuration belongs to the abstract register file
assigned to its decoded instruction point.

\paragraph{Proof.}
The initial seed is included in its point annotation.  Each concrete source
step is included in the corresponding abstract transfer just described.  The
post-fixpoint condition says that the resulting abstract file is included in
the annotation at its actual successor.  Induction on the number of source
steps proves the assertion.  The induction ranges over arbitrary natural
step counts.  All underlying objects are finite strings, natural numbers, and
finite lists; the argument is formalizable in ZFC.  The Python replay is a
mathematical checker, not a new kernel-checked theorem.

\subsection{Exact consequences}

The resulting annotation reaches 6,452 instruction points and excludes exactly
three.  It proves the following stronger zero-register facts:
\[
\begin{array}{c|l|l}
\text{index}&\text{instruction path}&\text{registers forced to zero}\\\hline
6314&\texttt{111010010110001110}&R_3,R_6,R_7\\
6429&\texttt{111111000}&R_1,R_3,R_5,R_6,R_7
\end{array}
\]
The instructions at these two points are respectively $\operatorname{CD}(6)$
and $\operatorname{CD}(1)$.  Their positive branches are therefore impossible
in the initialized source computation.  The three excluded points are
\[
\begin{array}{c|l|l}
\text{index}&\text{path}&\text{instruction}\\\hline
6315&\texttt{111010010110001111}&\operatorname{CC}(2)\\
6430&\texttt{1111110010}&\operatorname{CI}(6)\\
6431&\texttt{1111110011}&\operatorname{CC}(2).
\end{array}
\]
Thus the second and third excluded leaves form one wholly unreachable subtree
at prefix \texttt{111111001}.  A simpler three-class abstraction using zero,
positive even, and odd already proves the last two leaves unreachable and the
zero fact at point 6429.  Increasing the exact-value threshold from 16 to 64
produces the same three excluded points and two forced decrement guards; this
is a scoped observation, not a completeness result for source invariants.

\subsection{A replacement theorem for the initialized computation}

At path \texttt{111010010110001110}, replace $\operatorname{CD}(6)$ by any of
\[
\operatorname{CB}(1),\quad
\operatorname{CD}(3),\quad \operatorname{CD}(6),\quad \operatorname{CD}(7).
\]
At path \texttt{111111000}, replace $\operatorname{CD}(1)$ by any of
\[
\operatorname{CB}(1),\quad
\operatorname{CD}(1),\quad \operatorname{CD}(3),\quad
\operatorname{CD}(5),\quad \operatorname{CD}(6),\quad \operatorname{CD}(7).
\]
Finally replace the subtree at either or both prefixes
\[
\texttt{111010010110001111},\qquad \texttt{111111001}
\]
by any finite source subtrees, leaving their attachment prefixes and all other
source structure fixed.  No remaining reachable instruction path is changed.

\paragraph{Source replacement theorem.}
Every such replacement gives exactly the same sequence of register files and
source program-counter words as the original program from the displayed seed.
In particular its source halting behavior is identical.

\paragraph{Proof.}
Compare the two source computations by induction.  The original computation
never decodes a point inside either replaced subtree, by the invariant.  At a
retained forced-zero guard, every permitted replacement is a zero conditional
decrement or an explicit level-one break.  Each preserves the entire register
file and applies the identical level-one-break operation to the identical
program-counter word.  Every other decoded instruction is unchanged.
Consequently one old step and one new step have exactly the same successor
configuration.  The induction keeps the two trajectories equal and hence also
keeps the new trajectory away from the replaced subtrees.  This proves both
halting and nonhalting equivalence without any consistency assumption.

This theorem is specific to the initialized seed.  It makes no claim about
arbitrary register files at the affected instructions.  A compiler or binary
runtime proof that requires every syntactic instruction guard must either
retain that broader interface or explicitly incorporate this new invariant.
The inherited ZFC calibration can be composed with the source equivalence
once a complete binary simulation and a correct blank initializer for a new
table have been established.

\subsection{A checked contraction and its present limitation}

As an additional experiment, remove the entire zero loop at prefix
\texttt{11111100}, consisting of old points 6429--6431, and contract its
surviving sibling \texttt{11111101} into \texttt{1111110}.  The resulting source
has 6,452 leaves.  Every surviving instruction label and jump operand remains
unchanged, and all its successor indices agree with the old graph after
sending the removed zero loop to old point 6432.  The original computation
can make only one consecutive step inside the removed loop, so deleting that
step is a proper stuttering simulation and preserves infinite computations.

However, hash-consing of the source decision tree, treating each complete
instruction as an atomic leaf, changes its number of unique internal subtrees
from 201 to 202.  Therefore this contraction does not
supply a smaller binary machine.  The more useful optimization interface is
the preceding replacement theorem, which preserves all reachable path lengths
and permits the unreachable regions to be chosen for decoder sharing.  No
236-state table or other sub-237-state binary result is claimed by this
invariant package.  Initializer correctness must be checked independently for
any later literal decoder change, because the initializer shares states with
the running program.
